\documentclass[12pt]{article}
\usepackage[a3paper, margin=1in]{geometry}
\usepackage{amsmath, amssymb, amsthm, ulem, xcolor}

\begin{document}

\section*{Domácí úkol VIII}

Vypracoval: Daniel \textit{``Randál"} Ransdorf \hfill Podpis: \rule{4cm}{0.4pt}

\begin{enumerate}
  
  \item \textbf{Grarfy a jejich doplňky:}
  
    Grafy $G = (V,E)$ a $G' = (V',E')$ jsou izomorfní právě tehdy, když existuje bijekce
    $f : V \to V'$ taková, že pro všechny $v_1,v_2 \in V$ platí
    \[
    \{v_1,v_2\} \in E \;\Leftrightarrow\; \{f(v_1),f(v_2)\} \in E'.
    \]

    Z logické ekvivalence dostáváme
    \[
    \{v_1,v_2\} \in E \;\Leftrightarrow\; \{f(v_1),f(v_2)\} \in E'
    \quad\Longleftrightarrow\quad
    \{v_1,v_2\} \notin E \;\Leftrightarrow\; \{f(v_1),f(v_2)\} \notin E',
    \]
    což je právě
    \[
    \{v_1,v_2\} \in \overline{E}
    \;\Leftrightarrow\;
    \{f(v_1),f(v_2)\} \in \overline{E'}.
    \]

    Tedy stejné zobrazení $f$ je izomorfismus mezi $G$ a $G'$ právě tehdy, když je izomorfismem 
    mezi jejich doplňky $\overline{G}$ a $\overline{G'}$. 

  \item \textbf{Různé stupně:}
  
    Pro $n=1$ takový graf existuje, bude mít jen jeden vrchol se stupněm $0$, tedy platí, že všechny
    vrcholy mají jiný stupeň. Dokažme, že to pro $n\ge2$ nejde.

    \definecolor{suddenblue}{gray}{0.95}
    \colorlet{suddenblue}{cyan!15}
    Předpokládejme graf $G = (V,E)$ na \colorbox{suddenblue}{$n \geq 2$} vrcholech, jehož všechny vrcholy mají různé stupně.
    Platí 
    \[max \left\{deg(v) \;\middle|\; v \in V \right\} \leq n-1,\]
    protože vrchol může být maximálně spojen $n-1$ hranami se všemi ostatními $n-1$ vrcholy.
    Zobrazení
    \[
    \deg : V \to \{0,1,\dots,n-1\}
    \]
    je prosté kvůli různosti stupňů. Jelikož obě množiny zobrazení mají stejnou velikost $n$, $\deg$ je dokonce bijektivní.
    
    To znamená, že v grafu musí existovat vrchol $v_0$ se stupněm $n-1 \colorbox{suddenblue}{$\geq $}1$, který je spojený se všemi
    ostatními vrcholy. Potom každý vrchol má alespoň jednu hranu.
    Zároveň ale v grafu musí existovat vrchol $v_1$ se stupněm $0$, který
    tedy nemá žádnou napojenou hranu. To je spor.
    
    Proto pro $n \geq 2$ takový graf neexistuje, tudíž existuje jen pro $n=1$.

\end{enumerate}

\end{document}
